\documentclass[10pt,a4paper]{amsart}
\usepackage[margin=19mm]{geometry}
\usepackage{amsmath,amssymb,mathtools}
\usepackage[scr=rsfs,cal=euler]{mathalfa}
\usepackage{xfrac}
\usepackage[unicode,hidelinks,bookmarks=false]{hyperref}
\newcommand{\ie}{i.e.\ }
\newcommand{\cf}{cf.\ }
\newcommand{\resp}{respectively\ }
\newcommand{\Subsection}{\subsection}
\providecommand{\nobreakdash}{\nobreak\hbox{-}}
\setlength{\emergencystretch}{2em}
\multlinegap0pt
\usepackage{bm}
\usepackage{bbm}
\usepackage{dsfont}
\usepackage{xspace}
\usepackage{cancel}
\usepackage{mathtools}
\usepackage{amsthm}
\makeatletter
\@ifundefined{Thm}{\newtheorem{Thm}{Thm}}{}
\@ifundefined{Prop}{\newtheorem{Prop}[Thm]{Proposition}}{}
\makeatother
\def \p {{\partial}}
\def\P{\mathbb{P}} %{{\bf P}}
\newcommand{\connect}{\xleftrightarrow}
\title[One-arm exponents of high-dimensional percolation revisited]{One-arm exponents of high-dimensional percolation revisited}
\author{Diederik van Engelenburg, Christophe Garban, Romain Panis, Franco Severo}
\date{Source: arXiv:2510.21595v1 (CC BY 4.0)}
\begin{document}
\maketitle

\noindent\emph{Source and licence.} One-arm exponents of high-dimensional percolation revisited. Version arXiv:2510.21595v1 (CC BY 4.0), distributed under the Creative Commons Attribution 4.0 International licence, as recorded in the arXiv licence metadata for this exact version and shown on the abstract page https://arxiv.org/abs/2510.21595v1. Full rights notice: licenses/arxiv-2510.21595-CC-BY-4.0.txt.\par\smallskip

\noindent\textit{Editorial scope.} Counted unit: the Prop whose source label is \texttt{prop:equivalence of sharp length} in the article (source0.tex lines 475--480, source label \texttt{prop:equivalence of sharp length}), delivered together with the article's own complete original proof of it (source0.tex lines 481--508, one \texttt{proof} environment of 28 lines). No mathematics is written, repaired or re-derived: statement and proof are the article's own text line for line. Required context retained uncounted: eq:input L(p) (lines 367--369); eq:input Phi(H\_n) (lines 371--373); eq:BK (lines 414--416); eq:comparison psi phi (lines 443--445). Every definition, display and label of those lines is retained unchanged. Omitted from this package and not redistributed: 1--366, 370--370, 374--413, 417--442, 446--474, 509--933. Nothing inside a retained excerpt is omitted or altered: every retained body is delivered exactly as the article's lines are written, so no omission receipt of any kind accompanies this package. Citations: the retained excerpts carry no citation command at all, so no bibliography entry of the article is transcribed and no reference is redistributed. Reference adaptation: none was needed and none was made; every reference inside the delivered unit resolves inside it, so no unresolved reference or citation is produced. Printed slips kept literal: no printed word, symbol, hypothesis or number of the counted statement or of its proof was altered. Layout and class: the article's own class is \texttt{article} with packages including amsfonts, amsthm, amsmath, amssymb, amscd, mathrsfs, eucal, graphicx, epstopdf, pict2e, epic, xcolor, float, mathtools; the wrapper is the lane's pinned \texttt{amsart} editorial wrapper at 10pt on A4, which restates the theorem-like environments the retained text needs and adds the article's own macro definitions that the retained text needs (\texttt{\textbackslash P}, \texttt{\textbackslash connect}, \texttt{\textbackslash p}), with extra wrapper packages (bm, bbm, dsfont, xspace, cancel, mathtools). The delivered numbering is the unit's own. Assets: no retained span contains a figure environment, a graphics-inclusion command, a PDF-page inclusion, a table or a TikZ picture; no image file is copied into this package, the package declares no asset dependency and no image file is redistributed with it. Licence: version arXiv:2510.21595v1 of this article is distributed under the Creative Commons Attribution 4.0 International licence, as recorded in the arXiv licence metadata for this exact version and shown on the abstract page https://arxiv.org/abs/2510.21595v1. A full rights notice is delivered at licenses/arxiv-2510.21595-CC-BY-4.0.txt. Characters: every retained excerpt is pure ASCII, so the delivered engine, XeLaTeX with its default Latin Modern text font, reports no missing character.
	\begin{equation}\label{eq:input L(p)}
		c_1(p_c-p)^{-1/2}\leq L(p)\leq C_1(p_c-p)^{-1/2}.
	\end{equation}

	\begin{equation}\label{eq:input Phi(H_n)}
		\varphi_p(\mathbb H_n)\leq C_1\exp\left(-c_1(p_c-p)^{1/2}n\right).
	\end{equation}

\begin{equation}\tag{BK}\label{eq:BK}
	\P_p[A\circ B]\leq \P_p[A]\P_p[B],
\end{equation}

\begin{equation}\label{eq:comparison psi phi}
	p\psi_{p}(S)\leq \varphi_{p}(S) \lesssim \psi_{p}(S).
\end{equation}

\begin{Prop}\label{prop:equivalence of sharp length}
 Let $d>6$ and $L\geq L_0$. There exist $c,C>0$ such that, for every $p<p_c$,
\begin{equation}
	c(p_c-p)^{-1/2}\leq \tilde{L}(p)\leq C(p_c-p)^{-1/2}.
\end{equation}
\end{Prop}

\begin{proof} By \eqref{eq:input L(p)}, it suffices to show that $\tilde{L}(p)\asymp L(p)$. First, observe that (by definition) $\tilde L(p) \leq L(p)$. 
We now prove that $\tilde{L}(p)\gtrsim L(p)$. 

Let $S\subset \Lambda_{\tilde{L}(p)}$ be such that $\varphi_p(S)\leq 1/e^2$. Let $N= (K+1) (\tilde{L}(p)+L)$ for some large enough constant $K\geq 1$ to be chosen later. We will show that $L(p)\leq N$, which will conclude the proof. 
Let $x\in \partial \Lambda_N$ and observe that
\begin{equation}
\{0\connect{\Lambda_N\:} x\} \subset \bigcup_{\substack{u_1\in S\\ v_1\notin S\\ u_1\sim v_1}} \{0\connect{S\:} u_1\} \circ \{u_1v_1 \text{ is open}\} \circ \{v_1 \connect{\Lambda_N\:} x\}.
\end{equation}
 Using \eqref{eq:comparison psi phi} and \eqref{eq:BK}, we obtain 
\begin{equation}
\varphi_p(\Lambda_N)\lesssim  \psi_p(\Lambda_N) \leq \sum_{x\in\partial\Lambda_N}  \sum_{\substack{u_1\in S\\ v_1\notin S \\ u_1\sim v_1}} \tau_p^S(0,u_1)\,p\,  \tau_p^{\Lambda_N}(v_1,x).
\end{equation}
By iterating this expansion $K$ times, we get (with the convention $v_0=0$)
\begin{equation}\label{eq:iterated-phi}
	\varphi_p(\Lambda_N)\lesssim  \sum_{x\in\partial\Lambda_N}  
	\sum_{\substack{u_1\in S\\ v_1\notin S \\ u_1\sim v_1}} \tau_p^S(v_0,u_1)\,p
	\ldots\sum_{\substack{u_{K}\in (S+v_{K-1})\\ v_{K}\notin (S+v_{K-1}) \\ u_K\sim v_K}} \tau_p^S(v_{K-1},u_K)\,p\, \tau_p^{\Lambda_N}(v_K,x).\end{equation}
Observe that by construction, any $v_K$ as above must lie in $\Lambda_N$, and 
\begin{equation}\label{eq:bound_vK}
\sum_{x\in  \p \Lambda_N} \tau_p^{\Lambda_N}(v_K,x)=\psi_{p}(\Lambda_N-v_K)\leq \psi_{p_c}(\Lambda_N-v_K)   
\leq 2dp_c^{-1} \max_{n\geq 0}\varphi_{p_c}(\mathbb H_n)\lesssim 1,
\end{equation}
where in the second inequality we used \eqref{eq:comparison psi phi}, and in the last one \eqref{eq:input Phi(H_n)}. Plugging \eqref{eq:bound_vK} in \eqref{eq:iterated-phi} and using translation invariance, we obtain
\begin{equation}\label{}
\varphi_p(\Lambda_N)  \lesssim  \varphi_p(S)^K \leq e^{-2K}.
\end{equation}
If $K$ is large enough, we deduce that $\varphi_p(\Lambda_N)\leq 1/e^2$, so that $L(p)\leq N\lesssim \tilde{L}(p)$. This concludes the proof.
\end{proof}

\end{document}
